\begin{enumerate}
    \item[\textbf{Problem 3.1.}] On what basis did the Lilliputians conclude that Gulliver needed 1728 times as much food as they did?

    \item[\textbf{Problem 3.2.}] How would the Lilliputians' conclusion change if they had thought about the exchange of energy between a person (of any size) and the surrounding environment?

    \item[\textbf{Problem 3.3.}] How do the surface area and volume of a sphere scale? Why? (Hint: Analyze spheres of radii 1 and R.)

    \item[\textbf{Problem 3.4.}] Explain what would happen to an angle between two lines inscribed on a balloon as it was inflated to a radius R from a radius of 1.

    \item[\textbf{Problem 3.5.}] Confirm that eq. (3.2) does adequately portray the straight line drawn in Figure 3.2.

    \item[\textbf{Problem 3.6.}] Show how the equation $y = mx^b$ becomes a linear equation in a log-log plot.

    \item[\textbf{Problem 3.7.}] Write eq. (3.6) in a form suitable for a log-log plot.

    \item[\textbf{Problem 3.8.}] Confirm the dimensional relationship of eq. (3.10).

    \item[\textbf{Problem 3.9.}] Use dimensional analysis to confirm that power is equal to the time rate of change of kinetic energy.

    \item[\textbf{Problem 3.10.}] Confirm the dimensional relationship of eq. (3.12).

    \item[\textbf{Problem 3.11.}] Confirm that eq. (3.13) does show that $v$ is independent of $L$.

    \item[\textbf{Problem 3.12.}] What is the hearing range of an elephant? A whale? How do these ranges compare with those of humans?

    \item[\textbf{Problem 3.13.}] Confirm that the dimensions of eq. (3.15) are 1/T.

    \item[\textbf{Problem 3.14.}] Confirm that the dimensions of eq. (3.16) are 1/T.

    \item[\textbf{Problem 3.15.}] Confirm that the dimensions of eq. (3.17) are 1/T.

    \item[\textbf{Problem 3.16.}] Under what conditions is eq. (3.24) dimensionally consistent?

    \item[\textbf{Problem 3.17.}] Confirm that the stress of eq. (3.27) satisfies eq. (3.24).

    \item[\textbf{Problem 3.18.}] Confirm that eq. (3.28) is correct.

    \item[\textbf{Problem 3.19.}] Show that the deflection $w_{mp}$ of a beam with bending stiffness $EI$, length $L$, and under a concentrated load $P$ is governed by the two dimensionless groups in eq. (3.34).

    \item[\textbf{Problem 3.20.}] Why is the torque, $\tau$, in the apparatus of Figure 3.12 a constant?

    \item[\textbf{Problem 3.21.}] Expressed in terms of the wheel's geometric and gravitational properties, what is the magnitude of the torque in Problem 3.20?

    \item[\textbf{Problem 3.22.}] Confirm that eq. (3.30) is dimensionally correct.

    \item[\textbf{Problem 3.23.}] Confirm that eq. (3.31) is dimensionally correct.

    \item[\textbf{Problem 3.24.}] Calculate and confirm the estimated number of revolutions in the last column of Table 3.1.

    \item[\textbf{Problem 3.25.}] Confirm that eq. (3.39) is the correct representation of eq. (3.37) in terms of the five scaling factors of a simple beam.

    \item[\textbf{Problem 3.26.}] Formulate a hypothesis to explain why a wood pigeon and a buzzard appear to have such different proportions of $W_{fm} / W_b$ in Figure 3.2.

    \item[\textbf{Problem 3.27.}] Show that the equation that describes the log-log plot of Figure 3.7 can be found to be $h \simeq 1.23l^{0.68}$, where $h$ and $l$ are, respectively, the nave height and church length rendered dimensionless by dividing each by 1 ft.

    \item[\textbf{Problem 3.28.}] Using reasoning similar to that which brought us to eq. (3.13), show that the maximum speed at which animals can run is independent of size.

    \item[\textbf{Problem 3.29.}] The velocity of blood in the aorta is related to the difference in pressure between the heart and the arteries. Find the relationship between the velocity of the blood and the pressure difference. (Hint: Use the work-energy theorem as we did for bird hovering in Section 3.3.2.)

    \item[\textbf{Problem 3.30.}] The stilt, a small long-legged bird, was described in \textit{Gulliver's Travels} as weighing 4.5 ounces and having legs 8 inches long. A flamingo has a similar shape and weighs 4 pounds. Apply scaling arguments to show that its legs should be about 20 inches long (as they actually are!).

    \item[\textbf{Problem 3.31.}] Given that a robin weighs about 2 ounces, could we scale the length of its legs from the stilt data provided in Problem 3.30? Explain your answer.

    \item[\textbf{Problem 3.32.}] It was discovered that a certain cucumber had cells that divided when they grew to 1.5 times the volume of the "remaining" cells. Cells normally divide so that the ratio between their surface and their mass remains constant. Is the cucumber described as a "normal" cucumber?

    \item[\textbf{Problem 3.33.}] Find the range of values of the variable $x$ for which the approximation $\cosh(x/\lambda) \simeq \frac{1}{2}e^{x/\lambda}$ is acceptable, for scaling factors $\lambda = 1$ and $\lambda = 6$. Plot both functions for each of the two scaling factors.

    \item[\textbf{Problem 3.34.}] An experiment to determine the natural or fundamental period of oscillation of a simple spring-mass system (see Figure P3.34) is set up as follows. A spring of stiffness $k$ is fixed at one end and connected to a mass $m$ at its other, with the mass being able to move along an ideal, frictionless air track. The mass is displaced a distance $x_0$ from its initial resting position, after which it oscillates along the air track around that initial position. The time needed for a complete oscillation-that is, the period-is measured several times for several periods in succession, with the results being compared to the theoretical formula for the period, $T$: $T = 2\pi\sqrt{m/k}$. Assuming that $k$ and $m$ are known and that the stopwatch used to measure the period is accurate to $\pm 1\%$. What are the possible pitfalls that might prevent the successful experimental determination of $T$?

    \item[\textbf{Problem 3.35.}] When the structural elements called beams vibrate... \textit{(Note: This sentence is cut off/incomplete in the original text).}

    \item[\textbf{Problem 3.36.}] A steel beam of length of 20 cm is to be used to model a prototype timber beam whose span is 3.6 m. (a) Verify that the dimensionless group containing the load, the modulus, and the length is $P/EL^2$. (b) If the timber beam is to carry a load of 9000 N at a point 1.5 m from the left end, what load must be applied to the model to determine whether the prototype can carry its intended load? (Assume that the load-carrying capacity is the only behavior of interest here.)

\end{enumerate}
